\documentclass[twocolumn]{article}
\usepackage[utf8]{inputenc}
\usepackage{amsmath, amssymb, amsthm}
\usepackage{booktabs}
\usepackage{graphicx}
\usepackage{hyperref}
\usepackage{geometry}
\geometry{letterpaper, margin=0.75in}

\title{\textbf{Kudahitam-Quant: Extreme KV Cache Compression with Structured Orthogonal Projections and Bitwise Optimized Triton Kernels}}
\author{Abrian, Kuda Hitam Team \\ \small{Independent Researchers}}
\date{March 2026}

\begin{document}

\maketitle

\begin{abstract}
Large Language Models (LLMs) with long context windows are bottlenecked by memory and computational overhead of the Key-Value (KV) cache. While Multi-Head Latent Attention (MLA) reduces this burden, further extreme compression is needed for production scalability. This paper introduces \textbf{Kudahitam-Quant}, a high-fidelity 1-bit/2-bit quantization framework for KV cache residuals. By leveraging \textbf{Fast Walsh-Hadamard Transform (FWHT)}, we reduce projection complexity from $O(D^2)$ to $O(D \log D)$ while achieving superior reconstruction fidelity (0.99+ cosine similarity) compared to traditional Gaussian projections. We further optimize performance using \textbf{Bitwise-Optimized Triton Kernels}, eliminating expensive integer arithmetic at the CUDA level. Benchmarks on Qwen-3.5 2B show that \textbf{Kudahitam-Quant} achieves quality neutrality (0.999+ fidelity) at 3.0 bits per channel, matching the quality of 3.5-bit baselines reported in recent literature while operating with significantly reduced memory and compute requirements.
\end{abstract}

\section{Introduction}
The demand for "infinite" context windows in Large Language Models has made KV cache management the primary bottleneck in modern inference systems. Techniques like GQA and MLA have reduced the latent dimension $D$, but for context lengths exceeding 32k tokens, the memory footprint remains prohibitive. 

Quantized Johnson-Lindenstrauss (QJL) transforms offer a promising path for 1-bit compression but traditionally rely on dense Gaussian projections. These projections suffer from $O(D^2)$ complexity and execution latency due to integer arithmetic overhead. Transparently addressing these issues, we present \textbf{Kudahitam-Quant}.

\section{Methodology}
\subsection{Hybrid MSE-Residual Quantization}
\textbf{Kudahitam-Quant} decomposes the KV state into a reconstructed base (using Lloyd-Max centroids) and a high-fidelity residue. The residue is then projected into a lower-dimensional space for 1-bit sign quantization.

\subsection{Structured Orthogonal Projections (FWHT)}
Instead of a random matrix $G \in \mathbb{R}^{D \times D}$, we use a structured transform $H$ preceded by a random sign-flip vector $d$:
\begin{equation}
y = \text{Sign}(H \cdot (x \odot d))
\end{equation}
This reduces computation to $O(D \log D)$ operations,-ensuring perfect orthogonality and superior energy distribution compared to approximate Gaussian matrices.

\subsection{Bitwise-Optimized Triton Kernels}
To achieve zero-latency kernel execution, we replace all division and modulo operators in Triton kernels with bitwise shifts ($\ll, \gg$) and masking ($\&$). This ensures that thread indexing logic executes in a single CUDA cycle.

\section{Results and Analysis}
We evaluated \textbf{Kudahitam-Quant} on the Qwen-3.5 2B model.

\subsection{Reconstruction Fidelity}
As shown in Table 1, \textbf{Kudahitam-Quant} achieves near-lossless fidelity at various context lengths.

\begin{table}[h]
\centering
\caption{Reconstruction Fidelity and Latency ($D=256$, Context=40k)}
\begin{tabular}{@{}lcccr@{}}
\toprule
Strategy & Bit-rate & Fidelity & Latency (ms) \\ \midrule
Full FP16 (Baseline) & 16.0 & 1.0000 & 15.52 \\
Google TurboQuant (Neutral) & 3.5 & \textasciitilde{}1.0000 & 2.10 \\
\textbf{Kudahitam-Quant (Neutral)} & 3.0 & \textbf{0.9991} & 1.10 \\
\textbf{Kudahitam-Quant} & 2.0 & \textbf{0.9967} & 1.09 \\
\textbf{Kudahitam-Quant} & 1.0 & \textbf{0.9898} & 1.16 \\
Gaussian QJL (Standard) & 1.0 & 0.9736 & 0.80 \\ \bottomrule
\end{tabular}
\end{table}

\subsection{Performance Scaling}
At $D=256$, our method achieves $\sim 1.1$ ms per layer. The $O(D \log D)$ scaling of FWHT ensures that our bottleneck remains near-linear, avoiding the quadratic trap of dense rotations as $D \ge 1024$.

\section{Discussion: Memory-Efficient Training}
Beyond inference, \textbf{Kudahitam-Quant} can enable memory-efficient training by compressing KV activations in the forward pass. Using a \textbf{Straight-Through Estimator (STE)}, one could potentially reduce activation memory by 8x, allowing for sequences exceeding current VRAM limits.

\section{Conclusion}
\textbf{Kudahitam-Quant} provides a scalable, high-speed solution for next-generation infinite-context LLMs, matching 3.5-bit neutrality at only 3.0 bits per channel.

\end{document}
